\documentclass[12pt]{article}
\usepackage[T1]{fontenc}
\usepackage[utf8]{inputenc}
\usepackage{lmodern}
\usepackage{textcomp}
\usepackage{enumitem}
\usepackage{geometry}
\usepackage{graphicx}
\usepackage{amsmath, amssymb}
\geometry{margin=1in}
\begin{document}
\begin{center}
Chemistry - Undergraduate Level Assessment 5 \
Total Marks: 40 \
Answer all questions.
\end{center}

\section*{Multiple Choice (10 marks)}
\begin{enumerate}[label=\arabic*.]
\item When the Heisenberg uncertainty principle is applied to a quantum mechanical particle in the lowest energy level of a one-dimensional box, which of the following is true?
\begin{enumerate}[label=(\alph*)]
    \item Momentum is known exactly, but no information about position can be known.
    \item Position is known exactly, but no information about momentum can be known.
    \item No information about either position or momentum can be known.
    \item Neither position nor momentum can be known exactly.
\end{enumerate}
\item Which of the following statements most accurately explains why the 1 H spectrum of 12 CHCl 3 is a singlet?
\begin{enumerate}[label=(\alph*)]
    \item Both 35 Cl and 37 Cl have I = 0.
    \item The hydrogen atom undergoes rapid intermolecular exchange.
    \item The molecule is not rigid.
    \item Both 35 Cl and 37 Cl have electric quadrupole moments.
\end{enumerate}
\item Which one of the following statements is true:
\begin{enumerate}[label=(\alph*)]
    \item Protons and neutrons have orbital and spin angular momentum.
    \item Protons have orbital and spin angular momentum, neutrons have spin angular momentum.
    \item Protons and neutrons possess orbital angular momentum only.
    \item Protons and neutrons possess spin angular momentum only.
\end{enumerate}
\item Which of the following is classified as a conjugate acid-base pair?
\begin{enumerate}[label=(\alph*)]
    \item HCl / NaOH
    \item H$_{3}$O+ / H$_{2}$O
    \item O$_{2}$ / H$_{2}$O
    \item H+ / Cl-
\end{enumerate}
\item Which of the following statements about the lanthanide elements is NOT true?
\begin{enumerate}[label=(\alph*)]
    \item The most common oxidation state for the lanthanide elements is +3.
    \item Lanthanide complexes often have high coordination numbers (\textgreater{} 6).
    \item All of the lanthanide elements react with aqueous acid to liberate hydrogen.
    \item The atomic radii of the lanthanide elements increase across the period from La to Lu.
\end{enumerate}
\end{enumerate}

\section*{True / False (10 marks)}
\textit{Answer True or False}
\begin{enumerate}[label=\arabic*.]
\setcounter{enumi}{5}
\item True or False: CH 3 Cl has a 1 H resonance approximately 1.55 kHz downfield from TMS in a 12.0 T magnetic field.
\item True or False: A magnetic field of 1.2 T is sufficient to achieve the polarization of 2.5 x 10^-6 for 13 C at 298 K.
\item True or False: Unpaired electrons are essential for both paramagnetism and ferromagnetism in materials.
\item True or False: Carbonic acid is not generally considered a toxic pollutant in water systems.
\item True or False: Ammonia (NH 3) is the strongest base in liquid ammonia.
\end{enumerate}

\section*{Long Form Response (20 marks)}
\textit{Answer each question with a clear and complete explanation.}
\begin{enumerate}[label=\arabic*.]
\setcounter{enumi}{10}
\item Discuss the orbital angular momentum quantum number, l, of the electron that is most easily removed from ground-state aluminum during ionization. What factors influence this?
\item Discuss the factors that contribute to the difference in first ionization energy between argon and neon, and explain why argon's ionization energy is lower.
\end{enumerate}

\vfill
\noindent\textit{End of Paper}
\end{document}